\section{Introduction}
\label{sec:intro}

Continuous collision detection (CCD) asks a question that sounds simple and is
not. Take two configurations of a mesh, one at the start of a time step and one
at the end, with every vertex moving along a straight line in between. Does any
pair of primitives touch during the step, and if so, when? The answer gates the
step size in cloth, deformable and rigid-body simulation. The two ways of
getting it wrong are not symmetric.

If a method reports a collision that does not happen, the solver takes a shorter
step than it needed to. Work is wasted and nothing is broken. If a method misses
a collision, or reports a time of impact \emph{after} the true one, the solver
steps straight through the contact, and the configuration it lands in contains
an interpenetration. One family of solvers is built on a barrier
potential~\citep{li2020ipc,li2021codim,ferguson2021rigid}, and its whole
argument depends on the state remaining intersection-free. For those solvers the
result is worse than degraded: the barrier is infinite on the wrong side, and
the optimisation has nowhere to go.

\begin{definition}[Conservative]
\label{def:conservative}
A CCD implementation is \emph{conservative} if it reports a collision whenever
the geometries intersect during the step, and the time of impact it reports is
at or before the true earliest one.
\end{definition}

\Cref{def:conservative} is the specification this work implements, and it is
deliberately stronger than the usual one. Reporting the Boolean correctly is not
enough. A query can agree with an exact reference on whether a contact happens
and still return a time after it, which costs the solver as much as a missed
collision does. Every accuracy measurement in this paper is therefore
\emph{signed}, and a time of impact later than the exact root counts as a
failure of the same severity as a miss.

The difficulty is that conservativeness is a statement about real arithmetic,
while implementations run in floating point. \citet{wang2021tight} showed that
most published CCD implementations are not conservative once the rounding is
accounted for, and they gave the first one that provably is. That
implementation, TightInclusion, replaces interval arithmetic with a custom
inclusion function and a certified error bound. It is the reference we measure
against throughout, used exactly as released.

\paragraph{This work.} We present SCCD, a two-phase CCD pipeline for triangle
\emph{and} quad meshes, with host and device implementations of both phases. Its
conservativeness is structural and not empirical (\cref{sec:guarantee}), and we
assess it against TightInclusion on a dataset whose roots are known exactly. The
contribution is the parallel design, its measured performance and the accuracy
assessment, and not an implementation of a known algorithm on new hardware.

The first contribution is a broad phase with no tuning parameter and no sort: a
two-dimensional uniform cell list whose cell size comes from the data itself.
Its two queries have different shapes and are built differently. Vertex-face
bins a box into every cell its extent touches and attributes a pair to the cell
holding the corner of the overlap (\cref{sec:bp:cell}). Edge-edge is one list
against itself, which admits one entry per box at its minimum corner and a
forward walk in row-major cell order, pruned by a prefix maximum per row of
cells, so no bound on how many cells a box may span is needed
(\cref{sec:bp:self}).

The second is a host narrow phase that packs single-instruction,
multiple-data (SIMD) lanes across queries. TightInclusion's search is a
priority queue ordered by the time lower bound, and its answer is the first box
it accepts. We show in \cref{sec:np:order} that keeping the minimum over
\emph{all} accepted boxes and pruning against it gives the same answer under
\emph{any} traversal order. That licenses a lane to hold boxes from different
queries (\cref{alg:host}), which turns a search with per-query divergence into
a vector kernel. Sibling boxes then inherit four of their eight corner values
from their parent, halving the dominant cost.

The third is a device narrow phase built around redistribution. The difficulty of a CCD query is heavy-tailed: on one scene most
queries need a single box, while $208$ of them need more than a million
(\cref{tab:difficulty}). A large per-block work stack keeps such a query inside
the block that seeded it while its neighbours idle. We use a deliberately
\emph{small} shared stack that spills early into a double-buffered global
queue, which the next launch redistributes across every block
(\cref{alg:device}). On the worst scene this is worth two orders of magnitude
(\cref{tab:stackcap}).

The fourth is a signed accuracy assessment against a certified reference. It
covers $11{,}630{,}064$ query-mode checks against exact symbolic roots, on both
processors and all six scenes, of which $11{,}044{,}766$ place a reported time
of impact beside an exact root. SCCD reports no missed collision and no late
time of impact, and reproduces TightInclusion's median and worst-case error to
every printed digit while running faster (\cref{sec:results:ref}).

We are equally explicit about where the pipeline does not win, and why. The device narrow phase is slower than the host on four of six scenes, for a
structural reason we can name, and that is enough to put the whole pipeline
behind on one of them. The host cell list beats the device broad phase at every
mesh size we measured, and by more as the mesh grows. The mode that trades accuracy for speed is \emph{slower} on
half the host scenes. \Cref{sec:limitations} collects those results, together
with the threats to validity that bound what the measurements can support.
